\section{Logical scope and literal polynomial definitions}
\label{sec:foundation}
\subsection{What the consistency interpretation asserts}
Throughout, $\N=\{0,1,2,\ldots\}$ and $\Z$ denotes the unrestricted integers.
The domain is part of an equation, not a presentation detail. An equation over
$\N$ is not thereby an equation over $\Z$, and an unrestricted construction
requires its own sign argument or an explicit domain conversion.

Let $\operatorname{Ref}_{\mathrm{ZFC}}(r)$ be the fixed arithmetized predicate
saying that $r$ codes a finite syntactic ZFC refutation. For each selected
endpoint $P$ on the stated domain $D$, the claim is
\begin{equation}\label{eq:calibration}
 \exists\boldsymbol{x}\in D^n\;P(\boldsymbol{x})=0
 \quad\Longleftrightarrow\quad
 \exists r\in\N\;\operatorname{Ref}_{\mathrm{ZFC}}(r).
\end{equation}
The ordinary transformations proved here are composed with the identified
inherited calibration. On that proof chain, formalization in ZFC yields
\begin{equation}\label{eq:consistency}
 \mathrm{ZFC}\vdash
 \left[\forall\boldsymbol{x}\in D^n\;P(\boldsymbol{x})\ne0\right]
 \longleftrightarrow\Con(\mathrm{ZFC}).
\end{equation}
This report does not assert that either sentence in
\eqref{eq:consistency} is true. In particular it does not prove consistency,
decide solvability, establish a global minimum, or supply a new end-to-end
proof-assistant kernel check or encoded object-language ZFC derivation.

For the direct families, the retained ordinary logical interfaces are
\begin{quote}\small
\file{KEndpoint.c222_iff_zf_refutation}\\
\file{ConstructibleChoice.zf_refutation_iff_zfc_refutation}.
\end{quote}
The first relates the fixed numerical calculus to a finite ZF refutation;
the second is the retained ZF-to-ZFC refutation equivalence. The arithmetic
front implements that calculus using finite syntax constructors, reciprocal
CRT carriers, power or factorial certificates as appropriate, and range tests.
The selected scalar branch retains its separately identified native/compiler
and refutation calibration at one coherent coefficient tuple. Its domain
and calibration are not interchangeable with the unrestricted direct route.

\subsection{The modular structure of the proof}
The proof chain has four logically different layers.
\begin{enumerate}
\item The inherited calibration identifies a fixed finite arithmetic relation
with syntactic refutation. Its provenance and exact numerical payloads are
part of the source closure.
\item Unbounded arithmetic arguments prove both witness directions for the
certificate relations. These include Pell classification, CRT choices,
factorial fixed points, integer signs, and rank or interpolation lemmas.
\item Exact polynomial transformations replace one finite presentation by
another. Integral substitutions, restored monic coordinates, graph identities
and coefficient collection establish the claimed zero-set correspondence.
\item Finite independent checkers verify the actual emitted arrays and graphs,
not merely an abstract instance of a theorem.
\end{enumerate}
A hash certifies file identity, not a theorem; an exhaustive coefficient
check certifies an identity, not an unbounded Pell assertion. Each layer is
used only for what it establishes. The new sections give complete arguments
for this round's transformations; the retained source closure supplies
the inherited arithmetic and logical inputs at their stated boundary.

\subsection{Explicit finite arrays define the direct equations}
For a quartic endpoint, the ordered \texttt{variables} array in
\file{equation.json} fixes the independent coordinates
$x_0,\ldots,x_{n-1}$. A term encoded as $[c,[i_1,\ldots,i_d]]$ denotes
$c\prod_{j=1}^{d}x_{i_j}$; the empty index list denotes the constant $c$.
Repeated indices are retained and give ordinary powers. The $m$ residual
rows $R_0,\ldots,R_{m-1}$ have integer coefficients and degree at most two.
The equation is exactly
\begin{equation}\label{eq:sos}
 P(\boldsymbol{x})=\sum_{j=0}^{m-1}R_j(\boldsymbol{x})^2=0.
\end{equation}
Every coefficient and every coordinate appears in a finite supplied array.
The adjacent \file{expanded_polynomial.json} stores the fully collected
polynomial, including its constant term. There is no factorial, variable
exponent, division, binomial coefficient, bounded quantifier, or sequence
symbol among its algebraic primitives. Such expressions occur only in the
proof explaining its zeros.

For either domain $D\subseteq\Z$, equation \eqref{eq:sos} is equivalent to
$R_j(\boldsymbol{x})=0$ for every $j$, because each square is nonnegative.
Exact degree four is established by a collected nonzero quartic coefficient,
not merely by the upper bound $\deg R_j\le2$. All independent coordinates
are counted; named arithmetic abbreviations are not new unknowns.

The height-one endpoints are defined directly by their collected coefficient
arrays. Their graph certificates prove the displayed energy identity and
its zero-set correspondence. The independent checker reconstructs the
energy and compares every collected coefficient. The height claim therefore
concerns the final polynomial, not only its generating rows.

\subsection{Arithmetic graphs define the scalar equations}
The scalar graph leaves are the stated independent variables, ordinary
integer constants, or fixed authentic numerical ports. The selected ten-positive-integer branch has its own strictly positive
domain and complete coefficient tuple. Each internal node is binary addition, subtraction, or multiplication.
An acyclic graph with one designated output therefore denotes a unique
element of an ordinary polynomial ring over $\Z$. Registers abbreviate
polynomials; they do not introduce hidden existential variables. There is
no division or radical gate.

For the selected ten-positive branch the entire tuple
$\beta,K_0,M_C,M_F$ is the Start-only4127 tuple of
Section~\ref{sec:start4127-positive10}, with $\beta=2^{193969}-1$.
All four complete minimal unsigned numerical payloads are supplied and
reconstructed from the finite tile language and coefficient recipe. The
$M_F$ port is already once shifted by $\beta$; it must not be shifted again.
The earlier BFS943/radix67 positive tuple and the source34 six-port
integer tuple are different objects and are not mixed with this tuple.
No new eleven-integer endpoint is claimed in this report.

\subsection{What each measure counts}
For a nonzero collected polynomial $P=\sum_\alpha c_\alpha x^\alpha$, write
\[
 \deg P=\max_{c_\alpha\ne0}|\alpha|,\qquad
 |\supp P|=\#\{\alpha:c_\alpha\ne0\},\qquad
 \htpoly(P)=\max_\alpha|c_\alpha|.
\]
Support includes a nonzero constant coefficient. The number of residual
rows is distinct from support and from the number of operations. Scalar
cost counts live binary graph gates, including squares and multiplication
by fixed constants; copying a register and reading a leaf are free. Dead
nodes do not count. Unless an expanded scalar polynomial is materialized,
its support and total coefficient height are unmeasured, not zero.

A point dominates another only when it is no worse in every displayed
measure and strictly better in at least one, with domain and degree held
fixed where required. This is a finite comparison among specified
presentations. Height one, smaller dimension, sparse support and low graph
cost are different objectives, and no single global simplicity ordering is
claimed.
